\section{Introduction}

Light field depth estimation provides dense geometric representations from multi-view angular data, serving as an essential component in modern 3D video systems, autonomous navigation, and immersive virtual reality applications . The reliability of these downstream systems depends heavily on the precision of the estimated depth maps. However, real-world environments are inherently composed of heterogeneous materials, including specular, translucent, and scattering surfaces. These non-Lambertian reflectance properties violate the photo-consistency assumption across different viewpoints, presenting a substantial challenge for accurate depth recovery. As multi-view imaging systems are increasingly deployed in unconstrained urban and indoor environments, resolving depth ambiguities caused by complex material interactions has become an essential requirement for robust 3D scene understanding.

Existing depth estimation approaches, particularly those relying on Epipolar Plane Images (EPIs), predominantly assume Lambertian reflectance, where pixel intensity remains constant across viewpoints . This assumption fails substantially in the presence of non-Lambertian surfaces. Specular reflections shift dynamically with viewpoint changes, destroying the linear EPI structures, while volumetric scattering eliminates a single well-defined surface depth. The severity of this physical violation is evident in empirical evaluations. Standard EPI-based baselines yield a Mean Absolute Error (MAE) of 0.411 in non-Lambertian regions, substantially exceeding the acceptable threshold of 0.25. Even in Lambertian regions with complex textures, EPI slope ambiguity limits depth precision, resulting in an MAE of 0.387 against a required threshold of 0.16. Compounding this architectural limitation is the critical scarcity of non-Lambertian training data. Available datasets often contain only a minimal number of relevant scenes (e.g., 4 scenes), creating a data starvation bottleneck that prevents conventional deep networks from generalizing. Traditional frequency-domain analyses, such as the 2D Discrete Fourier Transform (DFT), become ineffective when applied to discrete, low-resolution angular samplings (e.g., $9 \times 9$ light fields), as they fail to capture reliable material signatures.

To address these physical and data-driven limitations, we propose a Unified Dual-Mask Physical Model for light field depth estimation in non-Lambertian environments. Instead of relying on ineffective frequency-domain assumptions, we introduce a three-layer angular signal decomposition physical model, formulated as $I(u,v) = DC + a \cdot u + b \cdot v + \epsilon(u,v)$. This formulation effectively decouples ambient illumination ($DC$), disparity/depth (linear terms $a, b$), and material/reflectance properties (residual term $\epsilon$). By extracting geometric angular features, such as angular gradients and correlation coherence, a Random Forest classifier accurately generates pixel-level masks to distinguish Lambertian from non-Lambertian regions. Building upon this decomposition, we construct a GeometricDualMask network architecture. Utilizing the derived disparity magnitude $P(x,y)$ and material proportion $M(x,y)$, the dual-mask mechanism performs adaptive feature routing and depth regression within a unified framework. To overcome data scarcity and domain imbalance, we implement a domain-balanced sampling and mixed training strategy, employing a WeightedRandomSampler to maintain gradient and exposure equilibrium across diverse datasets, such as the 170-scene UrbanLF-Syn.

The main contributions of this work are:
\textit{ We propose a three-layer angular signal decomposition physical model and a geometric feature-based detection mechanism that decouples illumination, depth, and material properties, enabling precise pixel-level identification of non-Lambertian regions without relying on ineffective frequency-domain analysis.
} We develop a GeometricDualMask network architecture integrated with a domain-balanced mixed training strategy, which utilizes disparity and material masks for adaptive feature routing, effectively mitigating feature confusion in heterogeneous material scenarios and overcoming the data scarcity of non-Lambertian samples.
We conduct extensive cross-dataset evaluations on HCInew, UrbanLF-Syn, and non-Lambertian datasets, accompanied by comprehensive ablation studies and computational cost analyses. The proposed method achieves an overall MAE of 0.133 and a mixed urban MAE of 0.081, demonstrating significant improvements in complex texture and mixed-material environments.

The remainder of this paper is organized as follows. Section II reviews related work in light field depth estimation and non-Lambertian surface modeling. Section III details the proposed three-layer physical decomposition and the GeometricDualMask architecture. Section IV presents the experimental setup, cross-dataset results, ablation studies, and computational efficiency analysis. Section V concludes the paper.

Traditional light field depth estimation methods predominantly rely on hand-crafted photo-consistency metrics and epipolar geometry . Methodologically, these approaches strictly assume Lambertian reflectance, positing that pixel intensities remain invariant across different viewpoints. This assumption collapses in the presence of non-Lambertian surfaces, such as specular or translucent materials, where viewpoint-dependent reflectance variations invalidate the matching cost computations. At the technical level, traditional algorithms struggle with matching ambiguities in textureless or repetitive regions and exhibit high sensitivity to sensor noise, limiting their applicability in unconstrained environments.

To mitigate the limitations of hand-crafted features, early deep learning-based methods employ convolutional neural networks to extract depth cues from Epipolar Plane Images (EPIs) or multi-view inputs . While these models improve performance on standard benchmarks, they implicitly inherit the Lambertian assumption during training. Methodologically, the network architectures are designed to detect linear EPI structures; when non-Lambertian reflections disrupt these structures, the networks lack a physical decoupling mechanism, leading to severe feature confusion between geometric disparity and material reflectance. Technically, even within Lambertian regions containing complex textures, these models suffer from EPI slope ambiguity, which restricts the precision of depth regression and yields elevated error rates.

Recent advancements introduce attention mechanisms, dual-branch architectures, and defocus-based auxiliary heads to process complex scenes . Despite these architectural enhancements, such methods lack explicit physical models to disentangle environmental illumination, geometric parallax, and material properties. Methodologically, processing heterogeneous materials through a uniform feature extraction pipeline prevents the adaptive routing of material-specific cues. At the technical level, these data-driven models heavily depend on large-scale, diverse training datasets. Given the extreme scarcity of non-Lambertian light field data, which is often restricted to a minimal number of training scenes, these complex architectures are prone to overfitting on minority domains. The absence of domain-balanced training strategies exacerbates gradient imbalances, degrading cross-domain generalization and resulting in substantial depth errors in mixed-material environments.

The aforementioned limitations indicate that purely data-driven or purely geometry-based paradigms are insufficient for heterogeneous light field depth estimation. Therefore, there is a pressing need for an approach that integrates a physical signal decomposition model to decouple reflectance from geometry, utilizes a unified dual-mask architecture for adaptive feature routing, and incorporates domain-balanced training to overcome data scarcity.

To address the physical decoupling deficiency and feature confusion in existing approaches, in this paper, we propose a Unified Dual-Mask Physical Model for non-Lambertian light field depth estimation. Instead of relying on implicit Lambertian assumptions or purely data-driven attention mechanisms, our method explicitly disentangles environmental illumination, geometric disparity, and material reflectance through a three-layer angular signal decomposition physical model. Specifically, we formulate the angular signal as $I(u,v) = DC + a\cdot u + b\cdot v + \varepsilon(u,v)$, which isolates the residual term $\varepsilon(u,v)$ to capture non-Lambertian reflectance properties. Building upon this decomposition, we construct a Geometric Dual-Mask network architecture that adaptively routes features through a Lambertian Epipolar Plane Image (EPI) branch and a non-Lambertian geometric branch. To handle the severe data imbalance across different reflectance domains, we incorporate a domain-balanced sampling strategy and a mask-weighted fusion mechanism, ensuring robust depth regression in complex mixed-material scenes.

The main contributions of this work are summarized as follows:

\begin{itemize}
 \item \textbf{Contribution 1}: We propose a three-layer angular signal decomposition mechanism combined with a geometric feature-based non-Lambertian detection module. By extracting angular gradients and correlation coherence instead of relying on ineffective 2D-DFT frequency analysis in low-resolution light fields, we train a random forest classifier with five core geometric features. This achieves pixel-level classification and dual-mask generation, resolving the physical coupling between depth disparity and material reflectance.
 \item \textbf{Contribution 2}: We develop a Geometric Dual-Mask network architecture featuring dual-branch adaptive feature routing and balanced projection. The Lambertian branch integrates continuous piecewise smooth depth priors via a Planar Regularized Sampling Operator (PRSO), while the non-Lambertian branch utilizes geometric features to process X-shaped EPI patterns and bimodal distributions. A mask-weighted fusion module with a unified projection dimension ($FUSE\_DIM=96$) eliminates channel imbalance, preventing feature confusion when processing heterogeneous materials.
 \item \textbf{Contribution 3}: We introduce a domain-balanced sampling and mixed training strategy utilizing a WeightedRandomSampler to maintain exposure and gradient equilibrium across multi-domain datasets, such as the 170-scene UrbanLF-Syn. This strategy mitigates the generalization degradation caused by the scarcity of non-Lambertian training data and optimizes the network convergence in mixed-material environments.
\end{itemize}
Extensive experiments on multiple light field datasets, including HCInew, UrbanLF-Syn, and specific non-Lambertian benchmarks, validate the effectiveness of the proposed method. The quantitative results demonstrate that our model achieves an overall Mean Absolute Error (MAE) of 0.133 and a mixed urban scene MAE of 0.081, yielding significant improvements in complex texture and mixed-material scenarios compared to baseline EPI networks and recent dual-branch architectures.

The remainder of this paper is organized as follows. Section II reviews related work on light field depth estimation and non-Lambertian surface processing. Section III details the three-layer angular signal decomposition and the Geometric Dual-Mask network architecture. Section IV describes the domain-balanced training strategy and implementation details. Section V presents the experimental setup, ablation studies, and comparative results. Finally, Section VI concludes the paper.